\begin{afterword}
\summaryhead

The post starts from a remark attributed to Einstein. Asked what he would do if Eddington's 1919 eclipse observations failed to match General Relativity, he replied that he would feel sorry for the good Lord, because the theory is correct. This looks like arrogance, since no one can know a theory is right before it is tested.

The author answers with the arithmetic of the previous day's post. To pick the correct hypothesis out of 100 million takes about 27 bits of evidence, and this holds for hunches as much as for conscious reasoning. So when the equations of General Relativity first came to Einstein, he must already have held enough evidence to single them out. It would be a coincidence for that evidence to be just enough to find the theory and not enough to be sure of it, and since the brain is an inefficient processor, he probably had far more than a perfect reasoner would need. Seen this way, the remark is not appalling, and the theory was correct.

\responsehead

The post's central idea has real content. In a large space of possibilities, proposing the right hypothesis is itself a sign of information, not luck, and a theory can deserve confidence before the decisive test. The idea is older than the post: Peirce made the counting argument in 1903, and drew a different lesson about confidence, that the guessing faculty is ``not strong enough to be oftener right than wrong.'' The problems are in the case the post applies the idea to, and in what its argument can show.

The post misdescribes the case. The remark, as its usual source reports it, was made to a student after Einstein had a first report of Eddington's result, about what he would have said had the result gone the other way. The post presents it as confidence ``in advance of experimental test.'' It also names none of the evidence Einstein is said to have held. The best-known piece, the theory's account of Mercury's orbit in November 1915, was an ordinary reason for confidence, and the post does not mention it.

The argument can be made only after the fact. It starts from the premise that Einstein got the equations right, and before a test that is the open question. Einstein's own record shows the difficulty. In March 1914 he wrote that he no longer doubted the correctness of the Entwurf theory, whatever the eclipse showed, and that theory was wrong. His confidence in the wrong theory and in the right one was stated in almost the same terms. The post's closing reminder that General Relativity ``was correct'' is the premise of its argument, restated at the end.

Two steps are weaker than they look. The post asks how likely it is that Einstein had ``exactly'' 29.5 bits, a single point, when the question is whether he had less than about 36, a range of some seven bits. And its last step moves from what ``a perfect Bayesian'' could conclude from Einstein's evidence to Einstein's own confidence, which the post does not show was based on that evidence.

In short: a sound and older point about how much evidence it takes to find the right hypothesis, applied to a remark made after the test, by an argument that can justify confidence only once the result is known.
\end{afterword}
